\documentclass[12pt]{article}
\usepackage[T1]{fontenc}
\usepackage[utf8]{inputenc}
\usepackage{lmodern}
\usepackage{textcomp}
\usepackage{enumitem}
\usepackage{geometry}
\geometry{margin=1in}
\begin{document}
\begin{center}
Answer Key - Geography - Undergraduate Level Assessment 23 \
\small (Answers are ordered exactly as in the question paper.)
\end{center}

\section*{Multiple Choice}
\begin{enumerate}[label=\arabic*.]
\item \textbf{Answer:} C
\\ \textit{Options:} A) 10\%; B) 20\%; C) 40\%; D) 80\%
\\ \textit{Note:} The correct answer is 40\% because, as of 2020, approximately 40\% of the population in Sub-Saharan Africa was classified as living in extreme poverty, defined by living on less than \$1.90 a day, highlighting the region's significant economic challenges.
\item \textbf{Answer:} C
\\ \textit{Options:} A) 10\%; B) 20\%; C) 40\%; D) 80\%
\\ \textit{Note:} The correct answer is 40\% because data from health organizations indicate that approximately 40\% of adults aged 18 years or older in the United States were classified as overweight based on body mass index (BMI) measurements in 2016.
\item \textbf{Answer:} B
\\ \textit{Options:} A) strongly supported by most studies; B) supported mainly by cross-section, not time-series studies; C) supported mainly by time-series, not cross-section studies; D) generally repudiated by empirical studies
\\ \textit{Note:} The inverted U hypothesis, which posits that inequality initially rises with economic development and later falls, has primarily been substantiated by cross-sectional studies that analyze data at a specific point in time, rather than longitudinal time-series studies that observe changes over time.
\item \textbf{Answer:} C
\\ \textit{Options:} A) \$150,000; B) \$75,000; C) \$35,000; D) \$15,000
\\ \textit{Note:} The correct answer of \$35,000 reflects the disparity in global income distribution, where this amount places an individual in the top one percent of earners worldwide, highlighting the contrast between average incomes in wealthier nations and those in developing regions.
\item \textbf{Answer:} D
\\ \textit{Options:} A) 50\%; B) 60\%; C) 70\%; D) 80\%
\\ \textit{Note:} The correct answer of 80\% reflects a significant consensus among Americans regarding the vital role of free media in a democratic society, indicating a strong societal value placed on the protection of press freedom from government interference.
\end{enumerate}

\section*{True / False}
\begin{enumerate}[label=\arabic*.]
\setcounter{enumi}{5}
\item \textbf{Answer:} False
\\ \textit{Options:} A) True; B) False
\\ \textit{Note:} The correct answer is: the piers at Weston-super-Mare, Clevedon, Portishead and Minehead were served by the paddle steamers of P and A Campbell
\item \textbf{Answer:} True
\\ \textit{Options:} A) True; B) False
\\ \textit{Note:} According to the passage: The Surgeon General
\item \textbf{Answer:} True
\\ \textit{Options:} A) True; B) False
\\ \textit{Note:} According to the passage: humans
\item \textbf{Answer:} True
\\ \textit{Options:} A) True; B) False
\\ \textit{Note:} According to the passage: the phrase is now printed on newly issued Washington, D.C. license plates
\item \textbf{Answer:} True
\\ \textit{Options:} A) True; B) False
\\ \textit{Note:} According to the passage: pink flamingos
\end{enumerate}

\section*{Long Form Response}
\begin{enumerate}[label=\arabic*.]
\setcounter{enumi}{10}
\item \textbf{Answer:} Mykonos
\\ \textit{Note:} Expected answer: Mykonos
\item \textbf{Answer:} hazardous conditions
\\ \textit{Note:} Expected answer: hazardous conditions
\end{enumerate}

\vfill
\noindent\textit{End of Answer Key}
\end{document}